% ═══════════════════════════════════════════════════════════════
% MUSTERLÖSUNG ML-03a: Las habitaciones
% Spanisch Klasse 7 - Sequenz 3: Mi casa
% ═══════════════════════════════════════════════════════════════
% Fachfarbe: #c41e3a (Karminrot)
% Lösungsfarbe: ROT
% Gesamtpunkte: 20
%
% WICHTIG: Identisches Layout wie AB, nur mit roten Lösungen
% ═══════════════════════════════════════════════════════════════

\documentclass[11pt, a4paper]{article}

% ═══════════════════════════════════════════════════════════════
% SW-MODUS TOGGLE
% ═══════════════════════════════════════════════════════════════
\newif\ifswmodus
\swmodusfalse

% ═══════════════════════════════════════════════════════════════
% PAKETE
% ═══════════════════════════════════════════════════════════════

\usepackage[utf8]{inputenc}
\usepackage[T1]{fontenc}
\usepackage[ngerman]{babel}
\usepackage{helvet}
\renewcommand{\familydefault}{\sfdefault}

\usepackage[margin=2cm]{geometry}
\usepackage{enumitem}
\usepackage{tabularx}
\usepackage{booktabs}
\usepackage{multirow}
\usepackage{tikz}

\usepackage{xcolor}
\usepackage{framed}
\usepackage{tcolorbox}
\tcbuselibrary{skins}

\usepackage{fancyhdr}
\usepackage{lastpage}
\usepackage{needspace}

% ═══════════════════════════════════════════════════════════════
% FARBEN
% ═══════════════════════════════════════════════════════════════

\definecolor{spanisch-color}{HTML}{c41e3a}
\definecolor{grau-color}{HTML}{888888}
\definecolor{hellgrau-color}{HTML}{f5f5f5}
\definecolor{orange-color}{HTML}{b35c00}
\definecolor{loesung-color}{HTML}{c62828}

\definecolor{sw-dark}{HTML}{000000}
\definecolor{sw-gray}{HTML}{666666}
\definecolor{sw-white}{HTML}{ffffff}

\ifswmodus
    \colorlet{spanisch}{sw-dark}
    \colorlet{grau}{sw-gray}
    \colorlet{hellgrau}{sw-white}
    \colorlet{orange}{sw-dark}
    \colorlet{loesung}{sw-dark}
\else
    \colorlet{spanisch}{spanisch-color}
    \colorlet{grau}{grau-color}
    \colorlet{hellgrau}{hellgrau-color}
    \colorlet{orange}{orange-color}
    \colorlet{loesung}{loesung-color}
\fi

\newcommand{\fachfarbe}{spanisch}

% ═══════════════════════════════════════════════════════════════
% VARIABLEN
% ═══════════════════════════════════════════════════════════════

\newcommand{\thema}{Las habitaciones}
\newcommand{\sequenz}{03}
\newcommand{\block}{a}
\newcommand{\fach}{Spanisch}
\newcommand{\klasse}{7}
\newcommand{\maxpunkte}{20}
\newcommand{\datum}{\today}

% ═══════════════════════════════════════════════════════════════
% KOPF- UND FUSSZEILE
% ═══════════════════════════════════════════════════════════════

\pagestyle{fancy}
\fancyhf{}
\renewcommand{\headrulewidth}{0.4pt}
\renewcommand{\footrulewidth}{0.4pt}

\lhead{\fach{} -- Klasse \klasse}
\chead{\textcolor{loesung}{\textbf{ML-\sequenz\block{} (MUSTERLÖSUNG)}}}
\rhead{\datum}

\lfoot{}
\cfoot{}
\rfoot{Seite \thepage\ von \pageref{LastPage}}

% ═══════════════════════════════════════════════════════════════
% BEFEHLE
% ═══════════════════════════════════════════════════════════════

\newcommand{\punktebox}[1]{\hfill\fbox{\small \_/#1}}
\newcommand{\loesung}[1]{\textcolor{loesung}{\textbf{#1}}}
\newcommand{\luecke}[1]{\rule{#1}{0.4pt}}

\newtcolorbox{merkkasten}[1][]{
    colback=hellgrau,
    colframe=\fachfarbe,
    colbacktitle=white,
    coltitle=\fachfarbe,
    boxrule=1pt,
    arc=0pt,
    left=8pt, right=8pt, top=8pt, bottom=8pt,
    fonttitle=\bfseries,
    attach boxed title to top left={yshift=-2mm, xshift=5mm},
    boxed title style={boxrule=0pt, colback=white},
    title=#1
}

\newtcolorbox{vokabelkasten}{
    colback=white,
    colframe=\fachfarbe,
    boxrule=0.5pt,
    arc=0pt,
    left=5pt, right=5pt, top=5pt, bottom=5pt
}

\newtcolorbox{hinweis}{
    colback=white,
    colframe=orange,
    colbacktitle=white,
    coltitle=orange,
    boxrule=1pt,
    arc=0pt,
    left=5pt, right=5pt, top=5pt, bottom=5pt,
    fonttitle=\bfseries,
    attach boxed title to top left={yshift=-2mm, xshift=5mm},
    boxed title style={boxrule=0pt, colback=white},
    title=Hinweis
}

\newenvironment{aufgabe}[1]{%
    \needspace{5\baselineskip}%
    \nopagebreak[4]%
    \vspace{10pt}
    \noindent\textbf{\textcolor{\fachfarbe}{Aufgabe #1}}
    \nopagebreak[4]%
    \vspace{5pt}
    \nopagebreak[4]%
    \begin{list}{}{\leftmargin=0pt}
    \item[]
}{%
    \end{list}
}

\newcommand{\es}[1]{\textcolor{\fachfarbe}{\textbf{#1}}}

% ═══════════════════════════════════════════════════════════════
% DOKUMENT
% ═══════════════════════════════════════════════════════════════

\begin{document}

% ═══════════════════════════════════════════════════════════════
% KOPFBEREICH
% ═══════════════════════════════════════════════════════════════

\begin{center}
    {\Large\bfseries\textcolor{\fachfarbe}{Arbeitsblatt: \thema}}\\[0.3em]
    {\large\textcolor{loesung}{\textbf{--- MUSTERLÖSUNG ---}}}
\end{center}

\vspace{5pt}

\noindent
\begin{tabularx}{\textwidth}{|X|X|X|}
\hline
\textbf{Name:} \textcolor{loesung}{Musterschüler/in} &
\textbf{Klasse:} \klasse &
\textbf{Datum:} \datum \\
\hline
\end{tabularx}

\vspace{15pt}

% ═══════════════════════════════════════════════════════════════
% MERKKASTEN: Las habitaciones (mit Lösungen)
% ═══════════════════════════════════════════════════════════════

\begin{merkkasten}[Merkkasten: Las habitaciones de la casa]
\textbf{Die Räume des Hauses:}

\vspace{0.5em}

\begin{tabular}{llll}
\es{la cocina} & = \loesung{die Küche} & \es{el salón} & = \loesung{das Wohnzimmer} \\[0.8em]
\es{el dormitorio} & = \loesung{das Schlafzimmer} & \es{el baño} & = \loesung{das Badezimmer} \\[0.8em]
\es{el comedor} & = \loesung{das Esszimmer} & \es{el jardín} & = \loesung{der Garten} \\[0.8em]
\es{el garaje} & = \loesung{die Garage} & \es{la terraza} & = \loesung{die Terrasse} \\
\end{tabular}

\vspace{0.8em}

\textbf{Genus-Regel:} Endung \textbf{-a} $\rightarrow$ meist \es{la} | Endung \textbf{-o, -ón, -or} $\rightarrow$ meist \es{el}

\end{merkkasten}

\vspace{10pt}

% ═══════════════════════════════════════════════════════════════
% AUFGABE 1: Zuordnung Bild-Wort (4P) - LÖSUNG
% ═══════════════════════════════════════════════════════════════

\begin{aufgabe}{1}
Ordne die spanischen Wörter den Bildern zu. Schreibe den Buchstaben in die Lücke. \punktebox{4}
\end{aufgabe}

\begin{vokabelkasten}
\textit{Wörter:} \es{la cocina} | \es{el salón} | \es{el dormitorio} | \es{el baño}
\end{vokabelkasten}

\vspace{0.8em}

\begin{center}
\begin{tabular}{cccc}

% Bild A: Küche (Herd-Symbol)
\begin{tikzpicture}[scale=0.8]
    \draw[thick] (0,0) rectangle (2,1.8);
    \draw (0.3,0.3) rectangle (0.8,0.8);
    \draw (1.2,0.3) rectangle (1.7,0.8);
    \draw (0.3,1) rectangle (0.8,1.5);
    \draw (1.2,1) rectangle (1.7,1.5);
    \node at (1,-0.4) {\textbf{A}};
\end{tikzpicture}
&
% Bild B: Wohnzimmer (Sofa-Symbol)
\begin{tikzpicture}[scale=0.8]
    \draw[thick] (0,0) -- (2,0) -- (2,0.8) -- (1.8,0.8) -- (1.8,1.2) -- (0.2,1.2) -- (0.2,0.8) -- (0,0.8) -- cycle;
    \draw (0.4,0.6) rectangle (1.6,1);
    \node at (1,-0.4) {\textbf{B}};
\end{tikzpicture}
&
% Bild C: Schlafzimmer (Bett-Symbol)
\begin{tikzpicture}[scale=0.8]
    \draw[thick] (0,0) rectangle (2,0.4);
    \draw[thick] (0,0.4) -- (0,1) -- (0.4,1) -- (0.4,0.4);
    \draw[thick] (2,0.4) -- (2,0.6) -- (1.6,0.6) -- (1.6,0.4);
    \draw (0.5,0.5) rectangle (1.8,0.9);
    \node at (1,-0.4) {\textbf{C}};
\end{tikzpicture}
&
% Bild D: Badezimmer (Badewannen-Symbol)
\begin{tikzpicture}[scale=0.8]
    \draw[thick] (0,0) -- (0,1) -- (0.2,1.2) -- (1.8,1.2) -- (2,1) -- (2,0);
    \draw (0.2,0) -- (0.2,0.8) -- (1.8,0.8) -- (1.8,0);
    \draw (0.4,1.2) -- (0.4,1.4);
    \node at (1,-0.4) {\textbf{D}};
\end{tikzpicture}
\\[1.5em]
\loesung{la cocina} & \loesung{el salón} & \loesung{el dormitorio} & \loesung{el baño} \\
\end{tabular}
\end{center}

\textit{\textcolor{grau}{Bewertung: 1 Punkt pro richtige Zuordnung}}

\vspace{15pt}

% ═══════════════════════════════════════════════════════════════
% AUFGABE 2: Artikel einsetzen el/la (4P) - LÖSUNG
% ═══════════════════════════════════════════════════════════════

\begin{aufgabe}{2}
Ergänze den richtigen Artikel: \es{el} oder \es{la}. \punktebox{4}
\end{aufgabe}

\noindent
\begin{tabular}{ll@{\hspace{3cm}}ll}
a) & \loesung{la} cocina & b) & \loesung{el} jardín \\[1em]
c) & \loesung{la} terraza & d) & \loesung{el} garaje \\
\end{tabular}

\vspace{0.5em}
\textit{\textcolor{grau}{Bewertung: 1 Punkt pro richtigen Artikel}}

\vspace{15pt}

% ═══════════════════════════════════════════════════════════════
% AUFGABE 3: Haus beschriften (6P) - LÖSUNG
% ═══════════════════════════════════════════════════════════════

\begin{aufgabe}{3}
Beschrifte den Hausplan mit den spanischen Wörtern (mit Artikel). \punktebox{6}
\end{aufgabe}

\begin{center}
\begin{tikzpicture}[scale=1]
    % Hausumriss
    \draw[very thick] (0,0) rectangle (8,5);
    % Dach
    \draw[very thick] (0,5) -- (4,6.5) -- (8,5);

    % Trennlinien für Räume
    \draw[thick] (4,0) -- (4,5);
    \draw[thick] (0,2.5) -- (4,2.5);
    \draw[thick] (4,2.5) -- (8,2.5);

    % Raum 1: oben links (Wohnzimmer)
    \node at (2,4.2) {\textbf{1.}};
    \node at (2,3.5) {\loesung{el salón}};

    % Raum 2: oben rechts (Küche)
    \node at (6,4.2) {\textbf{2.}};
    \node at (6,3.5) {\loesung{la cocina}};

    % Raum 3: unten links (Schlafzimmer)
    \node at (2,1.7) {\textbf{3.}};
    \node at (2,1) {\loesung{el dormitorio}};

    % Raum 4: unten rechts (Badezimmer)
    \node at (6,1.7) {\textbf{4.}};
    \node at (6,1) {\loesung{el baño}};

    % Garten (außen rechts)
    \draw[thick, dashed] (8,0) -- (10,0) -- (10,2.5) -- (8,2.5);
    \node at (9,1.7) {\textbf{5.}};
    \node at (9,1) {\loesung{el jardín}};

    % Esszimmer (Dachgeschoss-Bereich)
    \draw[thick] (2.5,5) -- (2.5,5.8) -- (5.5,5.8) -- (5.5,5);
    \node at (4,5.6) {\textbf{6.}};
    \node at (4,5.2) {\footnotesize\loesung{el comedor}};

\end{tikzpicture}
\end{center}

\textit{\textcolor{grau}{Bewertung: 1 Punkt pro richtige Beschriftung (inkl. korrektem Artikel)}}

\vspace{10pt}

% ═══════════════════════════════════════════════════════════════
% AUFGABE 4: Mini-Dialog ¿Dónde está...? (6P) - LÖSUNG
% ═══════════════════════════════════════════════════════════════

\begin{aufgabe}{4}
Schreibe einen Mini-Dialog. Verwende \es{¿Dónde está...?} und \es{Está en...} \punktebox{6}
\end{aufgabe}

\begin{vokabelkasten}
\textbf{Redemittel:}\\
\es{¿Dónde está tu madre/padre/hermano/hermana?} -- Wo ist deine Mutter/dein Vater/...?\\
\es{Está en la cocina / el salón / ...} -- Er/Sie ist in der Küche / im Wohnzimmer / ...
\end{vokabelkasten}

\vspace{0.8em}

\textbf{Beispiellösung:}

\vspace{0.5em}

\noindent
\textbf{A:} \es{¿Dónde está} \loesung{tu madre}\es{?}

\vspace{0.5em}

\noindent
\textbf{B:} \loesung{Está en la cocina.}

\vspace{0.8em}

\noindent
\textbf{A:} \es{¿Y} \loesung{tu padre}\es{?}

\vspace{0.5em}

\noindent
\textbf{B:} \loesung{Está en el jardín.}

\vspace{0.8em}

\noindent
\textbf{A:} \es{¿Dónde} \loesung{duermes tú}\es{?}

\vspace{0.5em}

\noindent
\textbf{B:} \loesung{Duermo en el dormitorio.}

\vspace{1em}

\textit{\textcolor{grau}{Bewertung: Je 2 Punkte pro sinnvollen Austausch (Frage + passende Antwort). Variationen akzeptiert.}}

% ═══════════════════════════════════════════════════════════════
% PUNKTEÜBERSICHT
% ═══════════════════════════════════════════════════════════════

\vspace{20pt}
\hrule
\vspace{10pt}

\begin{flushright}
\textbf{Punkte:} \loesung{\maxpunkte} / \maxpunkte
\end{flushright}

\end{document}
